%%% Основні відомості для заповнення%%%

\newcommand{\theWordDone}{Виконав}               % Виконав vs Виконала

\newcommand{\theStudent}{студент}                % студент vs студентка
\newcommand{\theSYear}{4}                        % Курс: 4
\newcommand{\theGroup}{ІО-25}                    % Група
\newcommand{\theStampCode}{ІАЛЦ.467200}          % Код у штампі: код КПІ та код пристрою 



\newcommand{\thesisAuthor}{Льоскін Іван Вадимович}             % Автор дипломного проекту
\newcommand{\thesisAuthorTask}{Льоскіну Івану Вадимовичу}       % Кому видане завдання для дипломного проектування
\newcommand{\thesisAuthorInitials}{Льоскін~І.~В.}              % Прізвище та ініціали автора дипломного проекту

\newcommand{\taskDate}{10~травня 2026 року}      % Дата видачі завдання
\newcommand{\completeDate}{5~червня \the\year\ року}  % Дата здачі закінченого проекту
\newcommand{\approvalDate}{5~червня 2026 року}  % Дата затвердення наказом по університету
\newcommand{\approvalNum}{1212-­с}   % Номер наказу затвердення теми

\newcommand{\thisisHD}           % Завідувач кафедри:  ім’я, прізвище
{\todo{Артем ВОЛОКИТА}}
\newcommand{\thisisHDInitials}   % Завідувач кафедри:  прізвище, ініціали
{Волокита~А.~М.}

\newcommand{\thesisTitle}        % Тема дипломного проєкту
{Кросплатформна клієнт-серверна система агрегації та обробки медико-біологічних даних з використанням хмарних технологій}
%{Інтерфейс користувача для класифікації рукописних цифр}

\newcommand{\thesisOsvitnyaPrograma}     % Освітня програма навчання
{\todo{Комп’ютерні системи та мережі}}

\newcommand{\thesisSpecialtyNumber}    % Номер спеціальності
{123}

\newcommand{\thesisSpecialtyTitle}     % Назва спеціальності
{Комп’ютерна інженерія}

\newcommand{\thesisDegree}             % Освiтньо-квалiфiкацiйний рiвень
{Бакалавр}


\newcommand{\thesisOrganization}       % Назва навчального закладу
{Національний технічний університет України}
\newcommand{\thesisOOrganization}
{"Київський політехнічний інститут імені Ігоря Сікорського"}

\newcommand{\thesisInstitute}           % Факультет
{Факультет інформатики та обчислювальної техніки}

\newcommand{\thesisDepartment}          % Кафедра
{Кафедра обчислювальної техніки}


\newcommand{\supervisorFio}            % Керівник проекту: прізвище, ім’я, по батькові
{Кобилюк Андрій Григорович}
\newcommand{\supervisorInitials}       % Керівник проекту:  прізвище, ініціали
{Кобилюк~А.~Г.}
\newcommand{\supervisorRegalia}        % Керівник проекту: посада, науковий ступінь, вчене звання
{\todo{посада, науковий ступінь}}

\newcommand{\visorFio}                 % Рецензент проекту: прізвище, ім’я, по батькові
{Сергієнко П.А.}
\newcommand{\visorRegalia}             % Рецензент проекту: посада, науковий ступінь, вчене звання
{ас. каф. СПСКС, д-р ф.}

\newcommand{\consultantFio}            % Консультант проекту: прізвище, ім’я, по батькові
{Нікольський С.С.}
\newcommand{\consultantRegalia}        % Консультант проекту: посада, науковий ступінь, вчене звання
{ас. каф. ОТ}
\newcommand{\consultantInitials}       % Консультант проекту: прізвище та ініціали
{Нікольський~С.С.}
\newcommand{\consultantPart}       % Консультант проекту: назва розділу
{Нормоконтроль}
                     

%-------- Вихiднi данi до проєкту -------------
\newcommand{\dataInput}
{
технічна документація, теоретичні та статистичні дані, патенти на винахід, наукові публікації.
}

%-------- Змiст розрахунково-пояснювальної записки -------------
\newcommand{\explanatoryNoteContents}
{
опис предметної області і напрямків дослідження, аналіз та характеристика об’єкту досліджень, аналіз існуючих рішень, опис архітектури системи.
}

%-------- Перелік графічного матеріалу -------------
\newcommand{\graphicMaterial}
{
принципова, структурна та  функціональна схеми системи.
}


%-------- Консультанти розділів проєкту -------------
\newcommand{\consultantsTab}
{
\begin{table}[H]
\begin{center}
\begin{tabular}{ |>{\centering}m{0.5cm}|m{5cm}|>{\centering}m{5.5cm}|>{\centering}m{2cm}|>{\centering}m{2cm}|}
\hline
\rowcolor{lightgray!20}
  &  &  & \multicolumn{2}{c |}{Підпис, дата}  \tabularnewline \cline{4-5}
  
\rowcolor{lightgray!20} 
\multirow{-1}{*}{№} & \multicolumn{1}{c|}{\multirow{-1}{*}{\raggedright Розділ}} &\multicolumn{1}{c|}{\multirow{-1}{*}{\raggedright Консультант}}  & завдання видав & завдання прийняв \tabularnewline
 
\rowcolor{lightgray!20}  
&  &   \multicolumn{1}{c|}{\multirow{-2}{*}{\raggedright\scriptsize(прізвище та ініціали)}}  &  &    \tabularnewline
\hline
   
\rownumber & \consultantPart  & \consultantInitials &  &  \TBstrut \tabularnewline  
\hline

%\rownumber & Назва розділу консультацій   & \consultantInitials &  &  \TBstrut \tabularnewline  % Поле для іншого консультанта
%\hline

\end{tabular}
\end{center}
\end{table}
}


%-------- Календарний план -------------
\newcommand{\calendarPlan}
{
\begin{table}[H]
\setcounter{magicrownumbers}{0}
\begin{center}
\begin{tabular}{|>{\centering}m{0.5cm}|>{\raggedright}m{8cm}|>{\centering}m{5cm}|>{\centering}m{2cm}|}
\hline
\rowcolor{lightgray!20} № & Назва етапу виконання дипломного проекту & Термін виконання & Примітка  \tabularnewline
\hline

\rownumber & Затвердження теми проєкту  & 06.04.2026-07.06.2026 &  \TBstrut \tabularnewline
\hline

\rownumber & Вивчення та аналіз завдання  & 02.05.2026-08.05.2026 &  \TBstrut \tabularnewline
\hline

\rownumber & Розробка архітектури та загальної структури системи  & 09.05.2026-15.05.2026 &  \TBstrut \tabularnewline
\hline

\rownumber & Програмна реалізація системи   & 16.05.2026-29.05.2026 &  \TBstrut \tabularnewline
\hline

\rownumber & Оформлення пояснювальної записки  & 30.05.2026-07.06.2026 &  \TBstrut \tabularnewline
\hline

\rownumber & Передзахист  & 06.06.2026 &  \TBstrut \tabularnewline
\hline

\rownumber & Захист  & 20.06.2026 &  \TBstrut \tabularnewline
\hline

%\rownumber &   &  &  \TBstrut \tabularnewline  % поле для нового етапу
%\hline
\end{tabular}
\end{center}
\end{table}
}


{}